\chapter{Vancomycin Level}

\section*{What Is This Test?}
The vancomycin level measures the serum concentration of vancomycin, a glycopeptide antibiotic used to treat serious Gram-positive bacterial infections, particularly methicillin-resistant Staphylococcus aureus (MRSA), coagulase-negative staphylococci, and ampicillin-resistant enterococci. Vancomycin inhibits bacterial cell wall synthesis by binding to the D-Ala-D-Ala terminus of the peptidoglycan pentapeptide precursor, preventing cross-linking of peptidoglycan chains and resulting in bacterial cell death. Vancomycin is administered intravenously (except for oral vancomycin for Clostridioides difficile colitis, which is not systemically absorbed). Vancomycin has a narrow therapeutic window: efficacy requires an area under the curve (AUC) to minimum inhibitory concentration (MIC) ratio (AUC/MIC) of 400--600 for optimal bactericidal activity; toxicity, primarily nephrotoxicity, increases with sustained trough concentrations >15--20 mcg/mL. The 2020 consensus guidelines from the American Society of Health-System Pharmacists (ASHP), the Infectious Diseases Society of America (IDSA), the Society of Infectious Diseases Pharmacists (SIDP), and the Pediatric Infectious Diseases Society (PIDS) recommend AUC-guided dosing (with Bayesian software) rather than trough-only monitoring for serious MRSA infections. Vancomycin is eliminated largely unchanged (80--90\%) by glomerular filtration.

\section*{Alternative Names}
Vancomycin Level; Vancomycin Trough; Vancomycin Serum Concentration; Vanc Level; Vanc TDM

\section*{Principle}
Vancomycin is measured by particle-enhanced turbidimetric inhibition immunoassay (PETINIA) or homogeneous enzyme immunoassay (EMIT) on automated clinical chemistry analyzers. In the PETINIA assay, vancomycin in the patient sample competes with a vancomycin-coated particle conjugate for binding to anti-vancomycin antibody. The rate of particle aggregation is inversely proportional to the vancomycin concentration. Results are reported in mcg/mL (mg/L).

\section*{History}
Vancomycin was discovered in 1953 at Eli Lilly from a soil sample containing Amycolatopsis orientalis (formerly Streptomyces orientalis) from the jungles of Borneo. The original formulation (``Mississippi Mud'') contained many impurities and caused significant nephrotoxicity and ototoxicity. Improvements in purification in the 1970s--1980s resulted in the current purified formulation. Vancomycin entered widespread use in the 1980s with the emergence of MRSA. Vancomycin-resistant Enterococcus (VRE) and vancomycin-intermediate S. aureus (VISA) and vancomycin-resistant S. aureus (VRSA) have emerged due to acquisition of the vanA/vanB gene clusters, which alter the peptidoglycan target from D-Ala-D-Ala to D-Ala-D-Lac, reducing vancomycin binding affinity 1,000-fold.

\section*{Why This Test Is Ordered}
Therapeutic drug monitoring to ensure vancomycin levels are within the target range for efficacy and to minimize nephrotoxicity. AUC/MIC ratio of 400--600 is the target pharmacokinetic/pharmacodynamic parameter. Monitoring is indicated for: serious MRSA infections (bacteremia, endocarditis, osteomyelitis, hospital-acquired pneumonia, meningitis), patients with renal impairment or changing renal function, patients receiving concurrent nephrotoxic drugs (aminoglycosides, amphotericin B, contrast dye), obesity (altered volume of distribution), prolonged courses (>3--5 days), and critically ill patients.

\section*{Primary vs Secondary Test}
This is a primary therapeutic drug monitoring test for vancomycin.

\section*{How This Test Is Performed}
Blood sample in a serum separator tube (SST, gold-top) or red-top tube. For trough-only monitoring: draw blood within 30 minutes before the next dose. For AUC-guided monitoring: typically requires two serum concentrations (peak 1--2 hours after infusion end + trough) with Bayesian estimation (or one steady-state trough for Bayesian modeling).

\section*{What Preparation Is Needed}
Accurate documentation of the dose and the time of the blood draw is essential. No fasting is required. The trough is drawn within 30 minutes before the next dose and only after steady state has been reached (typically after 3--4 doses in patients with normal renal function). For AUC-guided monitoring, two serum concentrations (a post-infusion level 1--2 hours after the end of the infusion plus a trough) or a single steady-state trough for Bayesian modeling are used; the exact infusion start and stop times and draw times must be recorded. A baseline serum creatinine is measured before the first dose.

\section*{Contraindications and Precautions}
There are no contraindications to the blood draw. Precautions: Levels drawn before steady state or at the wrong time relative to the dose are misleading and must not be used to adjust therapy. The 2020 ASHP/IDSA/SIDP/PIDS guidelines recommend AUC-guided dosing (AUC/MIC 400--600) rather than trough-only monitoring for serious MRSA infections; trough-only monitoring may misclassify exposure in patients with fluctuating renal function, obesity, or augmented renal clearance. Concurrent piperacillin-tazobactam administration increases the risk of vancomycin-associated acute kidney injury, and the timing of both drugs should be documented. The vancomycin MIC (by broth microdilution) should be known because the AUC target is calculated against it.

\section*{Risks}
Minimal risk from the venipuncture: slight pain or bruising at the site, and rarely hematoma or vasovagal syncope. Vancomycin-associated toxicities (nephrotoxicity, infusion-related reactions such as red man syndrome, and rare ototoxicity) are unrelated to the blood draw and are the reason the level is monitored.

\section*{What Happens After the Test}
Results are typically available within hours and are interpreted by the pharmacist or clinician against the dosing regimen, renal function, and the vancomycin MIC. For serious MRSA infection, the dose is adjusted to achieve an AUC/MIC of 400--600; when AUC monitoring is not available, trough targets of 10--15 mcg/mL (non-severe) or 15--20 mcg/mL (severe) infection are used. Serum creatinine is rechecked at least 2--3 times weekly, and a repeat level is obtained after dose adjustment once the new steady state is reached.

\section*{Understanding Results}

\subsection*{Below Therapeutic Range (Subtherapeutic)}
Trough <10 mcg/mL or AUC/MIC <400. Inadequate dosing, augmented renal clearance (common in young trauma patients, burns, pregnant women), or obesity with increased volume of distribution. Risk of treatment failure, emergence of resistant organisms, and reduced clinical efficacy.

\subsection*{Above Therapeutic Range (Potentially Toxic)}
Trough >15--20 mcg/mL or AUC/MIC >600. Increased risk of nephrotoxicity (acute kidney injury, rise in serum creatinine). Nephrotoxicity risk increases exponentially with trough >20 mcg/mL. Concurrent piperacillin-tazobactam administration increases the risk of vancomycin-associated AKI (synergistic nephrotoxicity). Ototoxicity is rare with current purified IV vancomycin formulations at standard therapeutic levels.

\section*{Normal Results}
2020 ASHP/IDSA/SIDP/PIDS guidelines (for serious MRSA infections): AUC/MIC 400--600 (assuming vancomycin MIC of 1 mcg/mL by broth microdilution, target AUC 400--600 mg*h/L). Trough-only monitoring (alternative when AUC is not available): 10--15 mcg/mL for non-severe infections; 15--20 mcg/mL for severe infections (bacteremia, endocarditis, osteomyelitis, CNS infections, hospital-acquired pneumonia). Intermittent dosing peak (post-distributional): 20--40 mcg/mL 1--2 hours after infusion.

\section*{Follow-up Tests}
Serum creatinine and BUN before initiation and at least 2--3 times weekly during therapy. Vancomycin AUC calculation via Bayesian software (BestDose, Insight Rx, DoseMeRx, or manual two-point method). Repeat vancomycin levels every 3--5 days if the patient is clinically stable with stable renal function. Audiometry if symptoms of ototoxicity (tinnitus, hearing loss) develop.

\section*{References}
\begin{enumerate}
  \item Rybak MJ, Le J, Lodise TP, et al. Therapeutic monitoring of vancomycin for serious methicillin-resistant Staphylococcus aureus infections: a revised consensus guideline and review by ASHP/IDSA/SIDP/PIDS. Am J Health Syst Pharm. 2020;77(11):835--864.
  \item Liu C, Bayer A, Cosgrove SE, et al. Clinical practice guidelines by the Infectious Diseases Society of America for the treatment of MRSA infections. Clin Infect Dis. 2011;52(3):e18--e55.
  \item Neely MN, Youn G, Jones B, et al. Are vancomycin trough concentrations adequate for optimal dosing? Antimicrob Agents Chemother. 2014;58(1):309--316.
  \item Moellering RC Jr. Vancomycin: a 50-year reassessment. Ann Intern Med. 2006;144(5):318--324.
  \item Winter ME. Basic Clinical Pharmacokinetics. 5th ed. Baltimore: Lippincott Williams \& Wilkins; 2010.
\end{enumerate}

\clearpage
